\chapter{Surviving extensions, monodromy, and dynamics of the quotient}
\label{chap:extensiones-dinamica-cociente}
\index{survival!global extension}
\index{monodromy!nonadic return}
\index{measure!emissions catalog}

The language of paths requires a decisive precision. An APP path,
a visible emission, an enriched TPK history, and a global section
are not four names for the same object. They constitute four consecutive levels
of a single construction. The path preserves positional order;
the emission preserves a finite reading; the history adds signature,
carry, orientation, boundary, and memory; the global section requires all
of its prefixes to survive simultaneously at every depth.

This hierarchy makes it possible to separate two questions that had remained intermingled. The first is logical: under what hypotheses a finite prefix truly belongs to a compatible infinite extension. The second is dynamical: what it means to transport a single history and what it means to transfer all its compatible continuations simultaneously. The distinction is essential in APP\@. A given trajectory carries one state to another; a path transfer carries a distribution of extensions to another distribution.

This section resolves the finite problem at depth six. The
published deterministic chronology has visible-phase return and cannot select
a unique probability. By contrast, transport biased by the nonadic phase,
when each active channel is interpreted as the normalized renewal of its entire
fibre of compatible continuations, has a uniform nine-step return
over the (468) emissions. The weights (1/3) then cease to be
an external choice: they follow from the three occurrences of each signature
(+1,-1,0) in the oriented nonadic cycle. Normalization within each fibre remains
a precise mathematical premise ---the uniform count of all compatible
paths--- and must not be conflated with the successor of a single
trajectory. Passage to the inverse limit over all depths preserves a
distinct statement, formulated at the end of the section.

\section{\(5\) levels of the notion of path}

Let \(\Gamma_{\rm APP}\) be the category of cardinal paths of the APP torus\@.
For each depth \(m\), let \(\mathcal H_m\) be the finite set of
admissible enriched histories
\[
 h_m=(v_0,e_1,v_1,\ldots,e_m,v_m),
\]
where each arrow \(e_j\) transports all fields declared by TPK\@. If
\(n\geq m\), the application
\[
 p_{n,m}:\mathcal H_n\longrightarrow\mathcal H_m
\]
removes the tail beyond depth \(m\). The positional shadow
\[
 \operatorname{sh}_m:\mathcal H_m\longrightarrow\Gamma_{\rm APP}
\]
forgets the memory fiber and preserves the underlying APP path.

\begin{definition}[Hierarchy of extensions]
\label{def:jerarquia-extensiones}
We distinguish the following objects.
\begin{enumerate}
  \item An \emph{APP path} is a morphism of
  \(\Gamma_{\rm APP}\).
  \item An \emph{observable route} is a reading of that morphism that
  jointly conserves the sum and product leaves.
  \item A \emph{TPK extension} is a lift of the observable route to
  a history \(h_m\in\mathcal H_m\), with its residue, quotient,
  signature, orientation, boundary, carry, and record data.
  \item A \emph{surviving extension} is a finite history that admits
  compatible prolongations to every horizon.
  \item A \emph{global section} is a compatible family
  \(x=(h_m)_{m\geq0}\) of surviving extensions.
\end{enumerate}
In particular, the route is the shadow of the extension; it is not an external object
that the extension subsequently chooses.
\end{definition}

For \(h\in\mathcal H_m\), its cylinder of prolongations is
\[
 Z_m(h)=
 \{x=(h_n)_n\in\varprojlim_n\mathcal H_n:x_m=h\}.
\]
In particular, every family of \(Z_m(h)\) also satisfies
\(p_{n,\ell}(h_n)=h_\ell\) for \(n\geq\ell\); collections of
incompatible truncations are not admitted.
Let us also define
\[
 \mathcal S_m^{(N)}=p_{N,m}(\mathcal H_N),
 \qquad
 \mathcal S_m^\infty=\bigcap_{N\geq m}\mathcal S_m^{(N)}.
\]

\begin{theorem}[Conditional extension survivor criterion]
\label{thm:criterio-extension-superviviente}
Let us fix \(h\in\mathcal H_m\) and, for \(N\geq m\), write
\[
 \mathcal P_N(h)=\{g\in\mathcal H_N:p_{N,m}(g)=h\}.
\]
Suppose that the truncations are coherent and that each
\(\mathcal P_N(h)\) is finite. The following are equivalent:
\begin{enumerate}
  \item \(\mathcal P_N(h)\ne\varnothing\) for every \(N\geq m\), that is,
  \(h\in\mathcal S_m^\infty\);
  \item the tree of finite extensions of \(h\) is finitely branching
  and has a nonempty level at every depth;
  \item \(Z_m(h)\ne\varnothing\).
\end{enumerate}
Consequently, whenever the surviving finite truncations are nonempty, the set
\[
 \Omega_{\rm H}^{\rm surv}
 =\varprojlim_m\mathcal S_m^\infty
\]
is exactly the space of surviving global sections. The theorem does not
assert that these truncations are nonempty for every APP input: that
property must be proved by means of the survival protocol.
\end{theorem}

\begin{proof}
The implication from the third condition to the first two is obtained by
truncation. Now assume the first condition and order the elements of
\(\bigsqcup_{N\geq m}\mathcal P_N(h)\) by the prefix relation. Each level
is finite and nonempty, and coherence of the truncations makes this union
a finitely branching tree with nodes at every depth; this yields the
second condition. Conversely, from the second condition König's lemma
provides an infinite branch. Its truncations form a compatible family
belonging to \(Z_m(h)\), which is the third condition. The same construction
identifies the indicated inverse limit wherever its finite truncations are
nonempty.
\end{proof}

The theorem resolves the possible discrepancy between “extendable to each
horizon separately” and “possessing a global prolongation”: finiteness of the
branches is the hypothesis that permits the choice of a single compatible chain. It does not
claim that all APP paths belong to
\(\operatorname{sh}(\Omega_{\rm H}^{\rm surv})\). That coverage requires proving
that the enriched closures do not empty the corresponding fibers.

\section{Phase return and monodromy relation}

The return of the phase--length quotient is formalized separately in
\cref{def:odometro-9-12,prop:proyeccion-fase-longitud}:
\[
(b,r)\xrightarrow{\;9\;}(b+1,r),
\qquad
(b,r)\xrightarrow{\;108\;}(b,r).
\]
These equalities belong to the odometer \(9\times12\). They do not constitute the
proof of \(F_{\rm obs}^{108}=I\), which uses the observable state and appears
below.

Let \(\phi\) be the phase coordinate with values in \(\mathbb Z/9\mathbb Z\).
A visible-phase return satisfies
\[
 \phi(h_{k+9})=\phi(h_k).
\]
The complete record may, however, satisfy
\[
 B(h_{k+9})\ne B(h_k).
\]
This inequality is not a failure of periodicity: it is the memory of the circuit.
An ordinary action of \(C_9\) on the complete state would require the
ninth power to be the identity and would erase precisely that information.

To type the return, let \(H\) be a hexad and \(P_2(H)\) its set of
fifteen pairs. Let \(\mathcal B_k\) be the record states that survive at phase
\(k\), let \(\mathcal T_k\subseteq\mathcal B_k\times\mathcal B_{k+1}\) be the published transport relation, and
suppose a fiber encoder is given
\[
 c_k:\mathcal B_k\longrightarrow P_2(H).
\]
Only after declaring those encoders is the observable relation defined.
\[
 R_k=(c_k\times c_{k+1})(\mathcal T_k)
 \subseteq P_2(H)\times P_2(H).
\]
The relational monodromy of a loop is
\[
 \mathcal M_k
 =R_{k+8}\circ R_{k+7}\circ\cdots\circ R_k.
\]

\begin{proposition}[Classification of the nonadic return given the encoder]
\label{prop:clasificacion-retorno-nonadico}
Three logically distinct levels are given.
\begin{enumerate}
  \item Without a proof of functionality, the return is the relation
  \(\mathcal M_k\).
  \item If each \(R_k\) is the graph of a map
  \(\mu_k:P_2(H)\to P_2(H)\), the transport is the skew product
  \[
   T(k,p)=(k+1,\mu_k(p)),
  \]
  and
  \[
   T^9(k,p)=(k,M_k(p)),
   \qquad
   M_k=\mu_{k+8}\circ\cdots\circ\mu_k.
  \]
  \item If the \(\mu_k\) are bijective, \(M_k\) is an automorphic holonomy.
  The transport descends to an ordinary action of
  \(C_9\) on the complete state if and only if \(M_k=I\) in all
  fibers.
\end{enumerate}
\end{proposition}

\begin{proof}
The first claim is the definition of image and relational composition once
\((c_k)_k\) have been fixed. Under
functionality, nine consecutive applications keep the first
coordinate fixed and compose their actions on the second. The composition of
bijections is a bijection. Finally, \(T^9=I\) is exactly equivalent to
\(M_k=I\) for each phase and each point of the fiber.
\end{proof}

\begin{corollary}[No return under cumulative transport]
\label{pf84:E02-29}
Let \(q:X\to V\), \(T:X\to X\), and \(x\in X\) be given. If
\[
 q(T^9x)=q(x),\qquad
 T^9x=\mathcal M_x\,x,\qquad
 \mathcal M_x\,x\ne x,
\]
then the projection returns and the state does not:
\[
 q(T^9x)=q(x),\qquad T^9x\ne x.
\]
The pointwise hypothesis \(\mathcal M_x\,x\ne x\) is indispensable; the global
inequality \(\mathcal M_x\ne I\) does not exclude \(x\) from being a fixed point.
\end{corollary}

\begin{proof}
The first equality is a hypothesis. The second and the pointwise separation give
\(T^9x=\mathcal M_x\,x\ne x\). The statement would be refuted by a case that
satisfies the three premises and had \(T^9x=x\); it does not authorize replacing the
third premise by \(\mathcal M_x\ne I\).
\end{proof}

\begin{theorem}[Functional sufficiency of the enriched tuple]
\label{pf84:E02-30}
Let
\[
 E:\mathscr X_{\rm TPK}^{\rm enr}\longrightarrow\mathscr S_E
\]
the tuple that conserves the visible block, new blocks and carry, lifting
state, cylinders, boundary signature, orientation, observable, phase, and
dodecaphasic register. Let
\[
 T:\mathscr X_{\rm TPK}^{\rm enr}\to
 \mathscr X_{\rm TPK}^{\rm enr},
 \qquad
 O:\mathscr X_{\rm TPK}^{\rm enr}\to Y.
\]
When maps exist
\[
 \widehat T:\mathscr S_E\to\mathscr X_{\rm TPK}^{\rm enr},
 \qquad
 \widehat O:\mathscr S_E\to Y
\]
such that
\[
 T=\widehat T\circ E,\qquad O=\widehat O\circ E,
\]
the equality \(E(x)=E(y)\) implies \(T(x)=T(y)\), \(O(x)=O(y)\) and, by
determinism, the same continuations and subsequent outputs.
\end{theorem}

\begin{proof}
By factorization,
\[
 T(x)=\widehat T(E(x))=\widehat T(E(y))=T(y),
 \qquad
 O(x)=\widehat O(E(x))=\widehat O(E(y))=O(y).
\]
The first equality permits iteration of the argument. This is a sufficient
condition, not an unproved equivalence. Two states with the same tuple and
different transitions or outputs, access to an absent coordinate, or
querying a target value are falsifiers. The visible five-element microfiber
of \cref{pf84:E02-22} controls only the nonfaithfulness of the
word \(w_6\); it does not prove this positive factorization.
\end{proof}

The published finite window contains nine phase returns with change of the
enriched state and a local selection from three continuations to one upon
increasing two horizons. This proves return with memory and refutes the
identification of the register with a memoryless cycle. Global extraction of
the nine encoders \(c_k\) and, through them, of the relations
\(R_k\), from all surviving histories is a distinct operation.
The preceding proposition classifies the return once the
register--pairs map is given; it does not derive that map from the enriched TPK state.

The concrete episode of the ninth phase, with its three continuations, the
block \(B_{11}\), the boundary \((502,662,638)\), and the integer inequalities
of cylinders, is proved only once in the
\cref{thm:seleccion-novena-fase}. In this residence only its
typed consequence is used:
\[
\#\operatorname{Surv}_{1}(B_9)=3,\qquad
\#\operatorname{Surv}_{2}(B_9)=1.
\]

The other two words can precede the same ternary tail only by changing
the deep Hensel extension. They share part of the visible residue,
but not the genealogy of the state.

\section{The catalog of 468 emissions as a product of channels}

The finite catalogue starts from the set
\(\Omega_{\rm seed}\) of \(104\,976\) microscopic presentations. Let
\(\mathcal U_6\) be the set of distinct sextuple emissions and let
\[
 F:\Omega_{\rm seed}\twoheadrightarrow\mathcal U_6.
\]
The exact census is
\[
 |\mathcal U_6|=468,
 \qquad
 \#\{U:|F^{-1}(U)|=144,432,1944\}=432,18,18.
\]
These counts belong to \(F\). Comparing them with depth-six TPK
histories would require an explicit map from the history domain
to \(\Omega_{\rm seed}\); no result in this section presupposes that
map or identifies catalogue presentations with enriched states.
The finite catalogue supplies, for every emission, an additive signature
\(s\) and a multiplicative signature \(e\). The joint map is a
bijection
\[
 \mathcal U_6\xrightarrow{\sim}
 \mathcal S_+\times\mathcal S_\times,
 \qquad
 |\mathcal S_+|=18,
 \quad
 |\mathcal S_\times|=26.
\]
Thus,
\[
 \boxed{468=18\cdot26}
\]
is a factorization of the object and not merely of its cardinality. The first
coordinate preserves how the emission sums; the second, how the
product accumulates. This separation is the explicit finite form of
sum--product incompatibility: both readings coexist in the state, but
each projection forgets the complementary signature.

The first coordinate admits an additional parametrization that will be useful when
comparing it with the nonadic boundary. If an additive seed has coordinates
\((i,j)\in(\mathbb Z/9\mathbb Z)^2\) and orientation
\(\varepsilon=+1\) for the east--south directions and \(-1\) for
north--west, its complete signature is determined bijectively by
\[
 \bigl(i+j\bmod9,\varepsilon\bigr)
 \in(\mathbb Z/9\mathbb Z)\times\{\pm1\}.
\]
Therefore,
\[
 \boxed{\mathcal S_+\simeq
 (\mathbb Z/9\mathbb Z)\times\{\pm1\}.}
\]
The product coordinate has no analogous uniformity: its \(26\) signatures
have fibers of seeds of sizes \(8\), \(24\), and \(108\), with histogram
\(24\times8+1\times24+1\times108=324\). This asymmetry is structural. The
factor \(18\) already contains phase and orientation; any subsequent descent
of the product factor to a class of pairs must preserve additional memory or
accept nonuniform fibers.

On \(\ell^2(\mathcal S_+\times\mathcal S_\times)\), let us define the
conditional expectations
\[
\begin{aligned}
 (E_+f)(s,e)&=\frac1{18}\sum_{s'\in\mathcal S_+}f(s',e),\\
 (E_\times f)(s,e)&=\frac1{26}\sum_{e'\in\mathcal S_\times}f(s,e').
\end{aligned}
\]
The first operator renews the additive signature and preserves the multiplicative
one; the second performs the dual operation. The third channel retains both.

\section{The unique path and path transfer}

Unlike the return of the odometer defined in
\cref{def:odometro-9-12,prop:proyeccion-fase-longitud}, the following proposition
proves the return of the observable operator \(F_{\rm obs}\).

The published observable law updates a single cursor at each step. In the
phases \(+1\), the additive cursor advances; in the phases \(-1\), the
multiplicative cursor advances; in the phases \(0\), both remain stationary. After the
steps \(27,54,81,108\), the two orientations are inverted. Denote this
deterministic evolution by \(F_{\rm obs}\).

\begin{proposition}[Return of the Observable Chronology]
\label{prop:retorno-fobs-identidad}
For every choice of the two initial positions and orientations, the
positions and orientations are restored after \(54\) steps. Upon
also restoring the phase, one obtains
\[
 F_{\rm obs}^{108}=I.
\]
Therefore, the Poincaré return of \(108\) steps induces the identity on
\(\mathcal U_6\) and does not select a unique stationary measure.
\end{proposition}

\begin{proof}
In every block of \(27\) steps, each cursor makes exactly nine advances in a fixed orientation. Since the position space is \(X_9=(\mathbb Z/9\mathbb Z)^2\), its displacement is \(9d=0\). At the end of the block, \(d\) is inverted. Two blocks restore position and orientation; four also restore the index in \(\Theta_{108}\). The induced return is consequently the identity. Every probability is stationary for the identity, so none is selected by that chronology.
\end{proof}

The problem arises even before the return if one attempts to forget the
route memory. The product signature
\[
 e_0=(5,5,5,5,5,5)
\]
has twenty-four representatives. After an advance in its own orientation,
those representatives separate, eight by eight, among
\[
 (5,5,5,7,1,7),\qquad
 (5,5,5,7,7,1),\qquad
 (5,5,5,1,7,7).
\]
Thus, the product advance does not descend to a map of the twenty-six
signatures: the representative that had been forgotten becomes relevant again. This
witness explains why an individual route cannot be turned, by mere
projection, into the refreshment of a complete fiber.

The APP reading is different. Instead of choosing a representative, it preserves all
compatible continuations. To formalize it in the finite catalogue,
consider the subalgebras of functions that depend only on one of the two
signatures. With respect to uniform counting, \(E_+\) and \(E_\times\) are the
conditional expectations that forget exactly the active coordinate. They are the
only Markov projectors that simultaneously satisfy:
\begin{enumerate}
  \item preservation of the coordinate inactive;
  \item complete forgetting of the active coordinate;
  \item preservation of the uniform counting state.
\end{enumerate}
Indeed, a row that has forgotten the source coordinate must be
constant on the \(18\), respectively \(26\), compatible destinations; the
normalization forces the values \(1/18\) and \(1/26\). This is the relative
canonicity of the renewal: no preference is postulated between branches that APP
declares equally possible.

Let \(C_9=\mathbb Z/9\mathbb Z\), with the phase word
\[
 \tau=(+1,-1,0,+1,-1,0,+1,-1,0).
\]
Write
\[
 P_{+1}=E_+,
 \qquad P_{-1}=E_\times,
 \qquad P_0=I.
\]

\begin{definition}[Phase-biased TPK transfer]
\label{def:transferencia-tpk-fase}
On \(\mathcal U_6\times C_9\) we define
\[
 \mathcal K_{\rm ph}\bigl((u,r),(v,r+1)\bigr)
 =P_{\tau(r)}(u,v),
\]
and set the remaining entries to zero. This definition is relative to the reader of
all compatible paths: an active phase uniformly renews the fiber
read by that phase, and the neutral phase retains the emission.
\end{definition}

\begin{theorem}[Uniform nonadic return]
\label{thm:retorno-nonadico-uniforme}
The kernel \(\mathcal K_{\rm ph}\) is doubly stochastic and irreducible,
has period exactly nine, and possesses a unique stationary state,
\[
 \pi_{\rm ph}(u,r)=\frac1{9\cdot468}.
\]
In each phase fiber, its return of nine steps is
\[
 \mathcal K_{\rm ph}^{\,9}\big|_r=E_+E_\times,
 \qquad
 (E_+E_\times)(u,v)=\frac1{468}.
\]
If the initial phase is uniformly distributed, the shadow of a step over
the emissions is
\[
 \overline K_{\rm ph}=\frac13(I+E_++E_\times).
\]
\end{theorem}

\begin{proof}
Each \(P_{\tau(r)}\) is doubly stochastic and the shift
\(r\mapsto r+1\) is a permutation; their skew product conserves both sums.
The three operators are idempotent and commute. Since each appears three
times in the nonadic word,
\[
 \prod_{r\in C_9}P_{\tau(r)}
 =E_+^3E_\times^3I^3=E_+E_\times.
\]
The first expectation forgets the additive signature and the second the multiplicative signature,
so their composition assigns to each of the \(18\cdot26=468\)
destinations the probability \(1/468\). Hence, in nine steps any two states of the same phase communicate; the
remaining steps communicate the phases, and the kernel is irreducible. Every return
must have length divisible by nine, whereas the diagonal of \(\mathcal K_{\rm ph}^9\) is positive;
the period is nine. Double stochasticity and irreducibility yield the unique
uniform state. Finally, averaging the nine phases counts three
copies of each channel and produces the final formula.
\end{proof}

Forgetting the phase is not a strong Markov quotient: from the same emission, two distinct phases apply \(I\), \(E_+\), or \(E_\times\), which have distinct rows. The operator \(\overline K_{\rm ph}\) is an averaged shadow in the stationary phase, not the chronological successor of a phaseless emission. This distinction preserves both the dynamics and the meaning of the measure.

\begin{definition}[Averaged kernel induced by transferrence]
\label{def:nucleo-tritico-refresco}
We call the phase shadow the averaged kernel of the catalog
\[
 \boxed{
 K_{\rm cat}=\frac13(I+E_++E_\times).
 }
\]
The weights are induced by the TPK calendar, which contains three occurrences
of each sign. The renewal operators are induced relative to the
uniform normalization of all compatible continuations of the active
fiber. Neither statement identifies \(K_{\rm cat}\) with
\(F_{\rm obs}\).
\end{definition}

\begin{proposition}[Convex family of stochastic kernels with a common uniform measure]
\label{prop:familia-refrescos-abc}
For \(a,b,c\geq0\), \(a+b+c=1\), the family
\[
 K_{a,b,c}=aI+bE_++cE_\times
\]
consists of symmetric, doubly stochastic kernels. If \(b,c>0\),
they are irreducible and, because they have positive diagonal, aperiodic. Their spectrum is
\[
 \operatorname{Spec}(K_{a,b,c})=
 \{1^{[1]},(a+c)^{[17]},(a+b)^{[25]},a^{[425]}\}.
\]
In particular, there exists a continuum of dynamics with the same factorization and the
same uniform measure. The choice
\(K_{\rm cat}=K_{1/3,1/3,1/3}\) is not characterized by the census alone. It is
selected jointly by the TPK phase word and the normalized transfer
of the \cref{def:transferencia-tpk-fase}.
\end{proposition}

\begin{proof}
The three summands are symmetric stochastic kernels; by convexity, the
same is true of \(K_{a,b,c}\), and symmetry converts row stochasticity
into column stochasticity. If \(b,c>0\), two steps make it possible to renew
both coordinates successively and connect any pair of states. Moreover,
\[
 K_{a,b,c}((s,e),(s,e))=a+\frac b{18}+\frac c{26}>0,
\]
so that the kernel is aperiodic. In the orthogonal decomposition
\[
 \mathbf1\oplus\ell_0^2(\mathcal S_+)
 \oplus\ell_0^2(\mathcal S_\times)
 \oplus\bigl(\ell_0^2(\mathcal S_+)\otimes
              \ell_0^2(\mathcal S_\times)\bigr),
\]
the pairs \((E_+,E_\times)\) act, respectively, as
\((1,1),(0,1),(1,0),(0,0)\). The four eigenvalues and the
multiplicities \(1,17,25,425\) are thereby obtained. For \(b,c>0\), irreducibility makes
the uniform stationary measure unique, which exhibits non-uniqueness even under
that requirement.
\end{proof}

\begin{theorem}[Exact spectrum of the balanced refreshment]
\label{thm:dinamica-cociente-468}
The kernel \(K_{\rm cat}\) is symmetric, doubly stochastic, irreducible and aperiodic. Its spectrum, with multiplicities, is
\[
 \boxed{
 \operatorname{Spec}(K_{\rm cat})
 =\left\{1^{[1]},
 \left(\frac23\right)^{[42]},
 \left(\frac13\right)^{[425]}\right\}.
 }
\]
Consequently, its spectral gap is (1/3) and its only stationary state is
\[
 \boxed{
 \tau_{468}(f)=\frac1{468}\sum_{U\in\mathcal U_6}f(U).
 }
\]
\end{theorem}

\begin{proof}
The operators \(E_+\) and \(E_\times\) are commuting orthogonal
averaging projectors. The rows and columns of each sum to one; the same
is true of their average. From \((s,e)\), \((s',e')\) is reached in two
renewals, and the contribution of \(I\) provides a self-edge. This
proves irreducibility and aperiodicity.

The space decomposes orthogonally into four sectors:
\[
 \mathbf1,
 \qquad
 \ell_0^2(\mathcal S_+),
 \qquad
 \ell_0^2(\mathcal S_\times),
 \qquad
 \ell_0^2(\mathcal S_+)\otimes
 \ell_0^2(\mathcal S_\times).
\]
Their dimensions are \(1,17,25,17\cdot25\). In the constant sector
the three summands act as the identity; in each one-coordinate sector
they act as (I+I+0); in the interaction sector, as (I+0+0).
Dividing by three gives the announced eigenvalues and multiplicities. A finite
irreducible and aperiodic kernel has a unique stationary state; since it is
doubly stochastic, that state is uniform.
\end{proof}

The theorem is exact relative to the factorization
\(\mathcal U_6\simeq\mathcal S_+\times\mathcal S_\times\), the TPK
calendar, and the uniform count of the compatible continuations of each fiber. The
family \(K_{a,b,c}\) proves that the dynamics do not follow from the census
alone. The \cref{prop:retorno-fobs-identidad} distinguishes the transfer of
all paths from the chronology of one of them.

\section{Compensated elevation constructed from the quotient}

Let us write \(m(U)=|F^{-1}(U)|\). From the kernel already chosen in the quotient
we define the kernel of presentations by
\[
 \boxed{
 K_{\rm seed}(x,y)
 =\frac{K_{\rm cat}(F(x),F(y))}{m(F(y))}
 }
\]
This kernel is stochastic. Summing over a target fiber yields
\[
 \sum_{F(y)=V}K_{\rm seed}(x,y)
 =K_{\rm cat}(F(x),V),
\]
independently of the representative \(x\). This identity shows that the
constructed lift projects exactly onto \(K_{\rm cat}\). It does not prove that
a preexisting microscopic dynamics of the enriched TPK state descends
to the catalog.

\begin{theorem}[Reversible elevation constructed from the quotient]
\label{thm:elevacion-reversible-cociente}
The probability
\[
 \widetilde\tau(x)
 =\frac1{468\,m(F(x))}
\]
is estacionaria and reversible for \(K_{\rm seed}\), and
\(F_\#\widetilde\tau=\tau_{468}\). If
\[
 J_F\delta_U
 =m(U)^{-1/2}\sum_{F(x)=U}\delta_x,
 \qquad
 D=\operatorname{diag}(m(U)),
\]
then
\[
 K_{\rm seed}J_F
 =J_F\bigl(D^{1/2}K_{\rm cat}D^{-1/2}\bigr).
\]
In particular, Hilbert-space intertwining without conjugation by \(D\)
occurs if and only if \(K_{\rm cat}\) commutes with the fiber
multiplicities.
\end{theorem}

\begin{proof}
Normalization follows because each of the 468 fibres receives total mass
(1/468). Using the symmetry of \(K_{\rm cat}\), for
\(F(x)=U\) and \(F(y)=V\) we have
\[
 \widetilde\tau(x)K_{\rm seed}(x,y)
 =\frac{K_{\rm cat}(U,V)}{468m(U)m(V)}
 =\widetilde\tau(y)K_{\rm seed}(y,x).
\]
This proves detailed balance, stationarity, and the identity of the
image measure.
Applying both sides of the last identity to a basis vector
\(\delta_V\) gives, on the fiber \(U\), the coefficient
\(K_{\rm cat}(U,V)/\sqrt{m(V)}\); that is exactly the coefficient of the
conjugate operator shown.
\end{proof}

\begin{theorem}[Phase Lift Faithful to Quiescence]
\label{thm:elevacion-fase-semillas}
Let \(X=\Omega_{\rm seed}\), \(m(u)=|F^{-1}(u)|\), and let us define
\[
\begin{aligned}
 L_0(x,y)&=\mathbf1_{\{x=y\}},\\
 L_+(x,y)&=\frac{E_+(F(x),F(y))}{m(F(y))},\\
 L_\times(x,y)&=\frac{E_\times(F(x),F(y))}{m(F(y))}.
\end{aligned}
\]
On \(X\times C_9\), let
\[
 \widehat{\mathcal K}_{\rm ph}\bigl((x,r),(y,r+1)\bigr)
 =L_{\tau(r)}(x,y).
\]
Then \((F,\mathrm{id}_{C_9})\) satisfies the property of
\emph{strong aggregability} (\emph{strong lumpability}) from
\(\widehat{\mathcal K}_{\rm ph}\) to \(\mathcal K_{\rm ph}\). The microscopic
kernel is irreducible, has period nine, and its unique stationary state
is
\[
 \widetilde\pi_{\rm ph}(x,r)
 =\frac1{9\cdot468\,m(F(x))}.
\]
Its nonadic return satisfies
\[
 \widehat{\mathcal K}_{\rm ph}^{\,9}
 \bigl((x,r),(y,r)\bigr)
 =\frac1{468\,m(F(y))}.
\]
\end{theorem}

\begin{proof}
The sum of \(L_+\), respectively \(L_\times\), over a target fiber
is \(E_+\), respectively \(E_\times\), and \(L_0\) projects to the identity.
This proves aggregability in each phase. The factors \(m(F(y))\) cancel
under composition: first the active signature is selected uniformly and then a
representative of the obtained emission. Thus, \(L_+\) and \(L_\times\) are
commuting projectors and
\[
 L_+L_\times(x,y)=\frac1{468\,m(F(y))}.
\]
The positivity of this entry connects all representatives in nine
steps. The phase forces the return lengths to be multiples of nine,
and the diagonal entry of the return is positive. The stationary formula is
verified by first summing the \(m(u)\) representatives of each emission; each
of the \(9\cdot468\) phase--emission classes receives the same mass.
\end{proof}

This lift improves the reading of the averaged kernel \(K_{\rm seed}\): in the neutral phase it preserves the microscopic representative, instead of silently remixing its fiber. Upon averaging and forgetting the phase, both constructions have the same shadow \(K_{\rm cat}\), but not the same hidden dynamics. The distinction is necessary in order to respect that the zero of TRIT is memory of quiescence.

\section{Minimal congruences stable under advance, half-turn and 4th of a turn: \(28\) and \(162\) memory classes}

The witness \(e_0=(5,5,5,5,5,5)\) shows that the product signature is too
coarse to transport a route. The exact missing minimal memory can be
computed. Let
\[
 \mathcal X_\times=X_9\times\mathscr D_{\rm card}
\]
the space of \(324\) product seeds. Let \(P\) denote the advance by one
cell in the stored orientation, \(H\) the half-turn of that
orientation, and \(Q\) the quarter-turn
\[
 N\longmapsto E\longmapsto S\longmapsto O\longmapsto N.
\]

\begin{theorem}[Forward and spin congruences]
\label{thm:congruencias-avance-giro}
The minimal congruence that refines the twenty-six product signatures and is stable
under \(P\) and \(H\) has \(28\) classes: twenty-seven have size eight, and one
has size \(108\). The only new splitting is
\[
 (5,5,5,5,5,5)\longrightarrow
 \mathcal L_0\sqcup\mathcal L_1\sqcup\mathcal L_2,
 \qquad |\mathcal L_j|=8.
\]
Its orbits under \(P,H\) have sizes
\[
 9,\quad9,\quad9,\quad1.
\]
By combining this refinement with the eighteen additive signatures one obtains
\[
 18\cdot28=504=468+36.
\]

The minimal congruence stable under \(P\) and \(Q\) has \(162\) classes of two
seeds. Each class is an orbit of the involution
\[
 J(i,j,d)=\bigl(7-i,\,7-j,\,\overline d\bigr)\pmod9.
\]
The operators induced by \(P,Q\) act transitively on the
\(162\) classes.
\end{theorem}

\begin{proof}
One begins with the partition by the twenty-six signatures and iteratively refines a class whenever two of its elements have images in distinct classes under any of the generators. The process terminates because \(\mathcal X_\times\) is finite. For \((P,H)\), the exhaustive census produces the histogram \(27\times8+1\times108=324\), the four indicated orbits, and three classes over the exceptional fiber of size twenty-four. For \((P,Q)\), it produces \(162\times2=324\). Direct verification proves that \(J\) preserves every signature, commutes with both generators, and pairs exactly the two elements of each class; the graph of classes generated by \(P,Q\) is connected.
\end{proof}

The excess \(36\) and the object \(R_{36}\) have distinct types. The boundary
at depth thirty-six is
\[
 R_{36}=
 \begin{pmatrix}
 222220\\021222\\102011
 \end{pmatrix},
\]
that is, three ordered blocks of six trits, one per channel. It is not a set of thirty-six states. The term \(504-468=36\), for its part, counts additional sheet memory: the half-turn that makes product transport functional adds exactly two new sheets to a single signature and, when repeated over the eighteen additive signatures, produces \(18(3-1)=36\). The coincidence of the integer (36) expresses two gradings of transport, boundary depth and sheet multiplicity, but does not authorize identification of their objects.

The quarter-turn \(Q\) was not hidden in \(F_{\rm obs}\) either. The
published chronology contains programmed advances and half-turns; incorporating
\(Q\) is an explicit enlargement of the route alphabet. Once it is incorporated,
a lazy symmetric random walk on \(P^{\pm1},Q^{\pm1}\) is irreducible and
aperiodic on the \(162\) classes, and its uniform measure is unique. This is the
exact mathematical form of “counting the turns”: the turn is a transport
coordinate and not a subsequent decimal correction.

\section{Mode and Orientation Cocycles}
\label{sec:cociclos-modo-orientacion}

The same rule makes it possible to define the operational counters without introducing
Dimensional Mechanics. Add to the state the last active mode
\(m\in\{\Sigma,\Pi\}\). A signature \(+1\) fixes \(m=\Sigma\), a signature
\(-1\) fixes \(m=\Pi\), and a signature \(0\) preserves the previous mode. For an
arrow \(a\), define the signed cocycle
\[
 \omega_{\rm sw}(a)=
 \begin{cases}
  +1,&\Pi\longrightarrow\Sigma,\\
  -1,&\Sigma\longrightarrow\Pi,\\
  0,&\text{the mode is preserved}.
 \end{cases}
\]
Its positive variation is \(s(a)=|\omega_{\rm sw}(a)|\). We further define
\(g(a)=1\) when the ordered pair of cardinal orientations changes and
\(g(a)=0\) when it is preserved. Keeping the mode and the terminal
orientation fixed at the endpoints, the sums
\[
 \Omega_{\rm sw}(\gamma)=\sum_{a\subset\gamma}\omega_{\rm sw}(a),
 \qquad
 S(\gamma)=\sum_{a\subset\gamma}s(a),
 \qquad
 G(\gamma)=\sum_{a\subset\gamma}g(a)
\]
are additive under concatenation.

In a nonadic closed loop, the sequence of modes is
\[
 (\Sigma,\Pi,\Pi,\Sigma,\Pi,\Pi,\Sigma,\Pi,\Pi).
\]
Therefore,
\[
 \Omega_{\rm sw}(C_9)=0,\qquad S(C_9)=6:
\]
there are three changes \(\Pi\to\Sigma\), three changes
\(\Sigma\to\Pi\), and three retentions. In the cycle of \(108\) steps,
\[
 \#(+1)=\#(-1)=\#(0)=36,\qquad
 \Omega_{\rm sw}=0,\qquad S=72,\qquad G=4.
\]
The final number counts the four simultaneous inversions programmed in
\(27,54,81,108\). These cocycles proceed directly from the TPK calendar\@.

\subsection{Quotient of modes and areal-volumetric correspondence}
\label{sec:cociente-modos-fav}

The calendar sign, mode label, and geometric sheet are three typed
objects. In particular, the phase sign is not identified with the
angular observable on digits. For the geometric correspondence, we fix the
typed HMT reader
\[
 \mathfrak g_{\rm av}(\Sigma)=\mathscr H_{\rm ar},
 \qquad
 \mathfrak g_{\rm av}(\Pi)=\mathscr H_{\rm vol}.
\]
The first assignment expresses the additive or boundary reading; the second, the
multiplicative or interior reading. The map \(\mathfrak g_{\rm av}\) is
part of the sheet typing. The following result proves that, once
that typing is fixed, TPK dynamics forces the incidence between them.

\begin{theorem}[Necessary descent of leaf incidence]
\label{thm:descenso-necesario-fav}
Let \(m_r\) be the mode obtained after applying phase \(r\) of the
primitive block \((+1,-1,0)\), with the rule
\[
 +1:\ m\mapsto\Sigma,\qquad
 -1:\ m\mapsto\Pi,\qquad
 0:\ m\mapsto m.
\]
The cyclic orbit of modes is \((\Sigma,\Pi,\Pi)\). If
\[
 N_3(i,j)
 :=
 \#\{r\in\mathbb Z/3\mathbb Z:(m_r,m_{r+1})=(i,j)\},
\]
then, in the order \((\Sigma,\Pi)\),
\[
 \boxed{
 N_3=
 \begin{pmatrix}0&1\\1&1\end{pmatrix}.}
\]
By repetition of the calendar,
\[
 N_9=3N_3,\qquad N_{27}=9N_3,\qquad N_{108}=36N_3.
\]
After applying \(\mathfrak g_{\rm av}\), the primitive incidence
matrix is, in the order
\((\mathscr H_{\rm ar},\mathscr H_{\rm vol})\),
\[
 \boxed{
 F_{\rm av}=
 \begin{pmatrix}0&1\\1&1\end{pmatrix}.}
\]
Thus, the two cross transitions, the unique volumetric self-retention, the
absence of areal self-retention and the unit primitive multiplicities
are consequences of the TPK mode quotient; they are not four added
postulates.
\end{theorem}

\begin{proof}
The three consecutive pairs, with cyclic closure, are
\[
 \Sigma\longrightarrow\Pi,\qquad
 \Pi\longrightarrow\Pi,\qquad
 \Pi\longrightarrow\Sigma.
\]
Each appears once, and the pair \(\Sigma\to\Sigma\) does not appear. This determines the four entries of \(N_3\). The calendars of lengths \(9,27,108\) respectively contain \(3,9,36\) copies of the primitive block, so their incidence matrices are the indicated multiples. The greatest common divisor of the nonzero entries of each long matrix removes only that global repetition and returns \(N_3\). Finally, \(\mathfrak g_{\rm av}\) transports \(\Sigma,\Pi\) to the areal and volumetric sheets without altering the incidence.
\end{proof}

The real hyperbolic realization that uses this matrix and proves its
Lorentz invariance is formulated in
\cref{sec:tres-generadores-concretos-trit}; there \(F_{\rm av}\) is used
as an antecedent, without repeating its derivation.

This theorem must be distinguished from a stronger Markov assertion.
Forgetting the phase is not a strong Markov quotient of
\(\mathcal K_{\rm ph}\): the visible mode does not determine which of the three
operators \(I,E_+,E_\times\) acts at the next elementary iteration. What descends
without an additional hypothesis is the edge multigraph of a TPK\@ block.

There is, however, a memory-one completion uniquely determined by that descent. Let \(X_{\rm av}\) be the smallest subshift of finite type over \(\{\mathscr H_{\rm ar},\mathscr H_{\rm vol}\}\) that contains all the words obtained by forgetting the phase and retains only consecutive pairs. Every one-step subshift with that property must admit exactly the three observed pairs; the smallest such subshift excludes the sole unobserved pair. Its adjacency matrix is therefore \(F_{\rm av}\).

\begin{corollary}[Parry Measure of the One-Memory Wrapper]
\label{cor:parry-envolvente-fav}
The subshift \(X_{\rm av}\) is primitive, has topological entropy
\(\log\varphi\), and possesses a unique measure of maximal entropy. Its stationary
vertex weights satisfy
\[
 \boxed{
 \frac{\mu_{\rm ar}}{\mu_{\rm vol}}=\varphi^{-2}.}
\]
\end{corollary}

\begin{proof}
The Perron root of \(F_{\rm av}\) is \(\varphi\), and a positive Perron
vector is \((1,\varphi)\). Since the matrix is symmetric, the Parry vertex
weights are proportional to the product of the left and right vectors,
that is, to \((1,\varphi^2)\).
\end{proof}

The Parry measure belongs to the memory-one path envelope. It is not
the image of \(\tau_{468}\), nor the temporal frequency of the periodic orbit
of modes. The latter is
\[
 \nu_{\rm cyc}(\mathscr H_{\rm ar})=\frac13,\qquad
 \nu_{\rm cyc}(\mathscr H_{\rm vol})=\frac23,
\]
whereas \(\tau_{468}\) is uniform over emissions, and the Parry measure
weights renewal histories of arbitrary length. The three measures
answer distinct questions. The irrational ratio
\(\varphi^{-2}\) is thus derived canonically from the path envelope,
without attributing it to a quotient of finite cardinalities.

\subsubsection{From chronological counting to the return marked by \texorpdfstring{\(\pi\)}{π}}
\label{sec:retorno-marcado-pi-phi2}

The two formulas
\[
 \left(\frac13,\frac23\right)
 \qquad\text{and}\qquad
 \frac{\mu_{\rm ar}}{\mu_{\rm vol}}=\varphi^{-2}
\]
are not approximations of one another. The first counts the three instants of the
primitive calendar \((\Sigma,\Pi,\Pi)\); the second weights all the
admissible histories of the memory-one envelope. Hence the chronological
frequency is rational, and the Parry ratio is irrational. In what follows,
we write
\[
 \frac{\pi_{\rm HMT}}{\varphi^2}
 =\pi_{\rm HMT}\varphi^{-2};
\]
this is not division by \(\varphi^{-2}\).

The appearance of \(\varphi^{-2}\) can be seen directly in the linear
dynamics. If \(F_n\) denotes the \(n\)-th Fibonacci number,
\[
 F_{\rm av}^{\,n}
 =
 \begin{pmatrix}
  F_{n-1}&F_n\\
  F_n&F_{n+1}
 \end{pmatrix}
 \qquad(n\geq1).
\]
The eigenvalues of \(F_{\rm av}\) are \(\varphi\) and \(-\varphi^{-1}\). The contracted branch reverses orientation in one iteration. Upon passing to the quadratic memory, namely \(\operatorname{Sym}^2(\mathbb R^2)\), that reversal is recorded before the sign disappears. In the basis \((x^2,2xy,y^2)\), with the images of the basis vectors represented by rows, the induced operator is
\[
 \operatorname{Sym}^2(F_{\rm av})=
 \begin{pmatrix}
  0&0&1\\
  0&1&2\\
  1&1&1
 \end{pmatrix},
 \qquad
 \chi(t)=(t+1)(t^2-3t+1).
\]
Its three eigenvalues are
\[
 \varphi^2,\qquad -1,\qquad \varphi^{-2}.
\]
Thus, \(\varphi^{-2}\) is the positive stable mode of second-order
memory. This step explains why the golden ratio appears here as a return
law and not as an added statistical mean.

The closure \(\pi\) belongs to another level of the same HMT state: it is the
coinductive output of the regional reader already constructed, not a new eigenvalue of
\(F_{\rm av}\). Let
\[
 P_{\rm ar}=
 \begin{pmatrix}1&0\\0&0\end{pmatrix},
 \qquad
 P_{\rm vol}=
 \begin{pmatrix}0&0\\0&1\end{pmatrix}.
\]
Among the positive operators that preserve both sheets, there is exactly
one whose volumetric normalization is one and whose areal mark is the generated
closure \(\pi_{\rm HMT}\):
\[
 D_\pi=\pi_{\rm HMT}P_{\rm ar}+P_{\rm vol}.
\]
Uniqueness is relative to these three explicit clauses: positivity,
diagonality with respect to the leaves, and the two normalizations. With
\(e_{\rm ar}=(1,0)^T\) and \(e_{\rm vol}=(0,1)^T\), the marked return of
length \(n\) is
\[
 \Theta_n
 :=
 \frac{\langle e_{\rm ar},
 D_\pi F_{\rm av}^{\,n}e_{\rm ar}\rangle}
 {\langle e_{\rm vol},
 F_{\rm av}^{\,n}e_{\rm vol}\rangle}
 =
 \pi_{\rm HMT}\frac{F_{n-1}}{F_{n+1}}.
\]
Therefore,
\[
 \boxed{\displaystyle
 \lim_{n\to\infty}\Theta_n
 =\frac{\pi_{\rm HMT}}{\varphi^2}.}
\]
The mark \(\pi_{\rm HMT}\) is applied only once, at the boundary return;
it is not reinjected at every transition. Finite combinatorics produces the
stable mode \(\varphi^{-2}\), the coinductive reader produces the closure
\(\pi_{\rm HMT}\), and \(D_\pi\) composes both data without reversing their
causal order.

The native length \(108\) provides an exact witness to that convergence:
\[
 (F_{\rm av}^{108})_{\rm ar,ar}
 =F_{107}=10284720757613717413913,
\]
\[
 (F_{\rm av}^{108})_{\rm vol,vol}
 =F_{109}=26925748508234281076009.
\]
Consequently,
\[
 \left|
 \frac{F_{107}}{F_{109}}-\varphi^{-2}
 \right|
 =
 6.1684979916614407936\ldots\times10^{-46},
\]
and, after the only closing marker,
\[
 \left|
 \pi_{\rm HMT}\frac{F_{107}}{F_{109}}
 -\frac{\pi_{\rm HMT}}{\varphi^2}
 \right|
 =
 1.9378907974286976068\ldots\times10^{-45}.
\]
These errors measure the finite approximation of the return; they are not adjustments of
\(\pi\) or of \(\varphi\).

\paragraph{Cylindrical reader of the ratio.}
Let \(I_n^\pi=[\underline\pi_n,\overline\pi_n]\) and
\(I_n^\varphi=[\underline\varphi_n,\overline\varphi_n]\) be the positive,
nested cylinders produced by the two HMT readers\@. The interval
operation
\[
 \mathcal Q(I,J)
 :=
 \left[
 \frac{\inf I}{(\sup J)^2},
 \frac{\sup I}{(\inf J)^2}
 \right]
\]
is the smallest interval containing \(x/y^2\) for every
\((x,y)\in I\times J\). It preserves nesting: if the input cylinders
are refined, then \(\mathcal Q(I_n^\pi,I_n^\varphi)\) is also refined. Since their
diameters tend to zero and the limit of \(I_n^\varphi\) is positive,
\[
 \bigcap_n\mathcal Q(I_n^\pi,I_n^\varphi)
 =
 \left\{\frac{\pi_{\rm HMT}}{\varphi_{\rm HMT}^2}\right\}.
\]
With the twelve already published base-thousand blocks, the quotient interval forces
thirty-five decimal digits:
\[
 \frac{\pi_{\rm HMT}}{\varphi_{\rm HMT}^2}
 =
 1.19998161486432666111577775331680692\ldots.
\]
This construction preserves the provenance of each digit: the numerator
comes from the closure cylinder and the denominator from the autoscale.

\paragraph{Algebraic separation between areal reading, volumetric reading and their respective codomains}
The complete ratio cannot arise from the rational algebra of
\(F_{\rm av}\) alone. Indeed, its rational spectral functions belong to
\(\mathbb Q(\sqrt5)\); \(\varphi^{-2}\) is found there, but not the
transcendental number \(\pi\). Consequently, any derivation purporting to
obtain \(\pi/\varphi^2\) solely from the matrix \(2\times2\) would be
ill typed. The transcendental datum must derive from the coinductive closure
reader, or from an equivalent boundary cocycle proved
independently. This negative control fixes the division of labor between
the finite correspondence and the number reader.

\subsubsection{Exchange involution of the quotient}
\label{sec:espejo-pi-phi-electron-hmt}

The correspondence admits two exact orientations of the same law. If
\[
 \mathscr R(x,y)=\frac{x}{y^2},
 \qquad
 \mathsf J(x,y)=(y,x),
\]
then \(\mathsf J^2=I\) and
\[
 \mathscr R(\pi,\varphi)=\frac{\pi}{\varphi^2},
 \qquad
 (\mathscr R\circ\mathsf J)(\pi,\varphi)
 =\frac{\varphi}{\pi^2}.
\]
The identities are also verified.
\[
 \mathscr R(\pi,\varphi)\mathscr R(\varphi,\pi)
 =\frac1{\pi\varphi},
 \qquad
 \frac{\mathscr R(\pi,\varphi)}
 {\mathscr R(\varphi,\pi)}
 =\left(\frac{\pi}{\varphi}\right)^3.
\]
Writing
\(\mathcal B_{\rm av}^{(\pi)}=\mathscr R(\pi,\varphi)\) also yields
the exact reparametrization
\[
 \frac{\varphi}{\pi^2}
 =
 \frac1{\varphi^3
 \bigl(\mathcal B_{\rm av}^{(\pi)}\bigr)^2},
 \qquad
 e^{\varphi/\pi^2}
 =
 \exp\!\left[
 \frac1{\varphi^3
 \bigl(\mathcal B_{\rm av}^{(\pi)}\bigr)^2}
 \right].
\]
Thus, the change of orientation between closure and growth transports the
boundary ratio to its exchanged reading:
\[
 \frac{\pi}{\varphi^2}
 \xleftrightarrow{\ \mathsf J\ }
 \frac{\varphi}{\pi^2}.
\]
The first orientation reads the marked areal closure over quadratic
volumetric memory; the second reads interior growth over
quadratic closure. This correspondence is a theorem of the HMT reader. Dimensional Mechanics subsequently applies the second orientation within the
electronic construction, without converting the electron into an antecedent of
\(\pi\), \(\varphi\), or of the correspondence.

\subsubsection{Ambivalence Principle of the areal and volumetric readers}
\label{sec:ambivalencia-grav-inercial-pi-phi2}

The Ambivalence Principle asserts that the areal reading and the
volumetric reading are two coordinated, but not interchangeable, functionals of one
and the same state with memory. Their coincidence on the diagonal is
\[
 \mathcal L_{\rm ar}^{\rm diag}(x)
 =\mathcal L_{\rm vol}^{\rm diag}(x).
\]
The diagonal equality does not erase the ambivalence outside of it.

On the areal--volumetric boundary, both readers factor through the same
positive memory amplitude \(\mathcal M(x)\):
\[
 \mathcal L_{\rm vol}(x)=\mu_{\rm vol}\mathcal M(x),
 \qquad
 \mathcal L_{\rm ar}(x)=
 \pi_{\rm HMT}\mu_{\rm ar}\mathcal M(x).
\]
Then the Parry theorem and the return mark give
\[
 \boxed{\displaystyle
 \frac{\mathcal L_{\rm ar}(x)}
 {\mathcal L_{\rm vol}(x)}
 =\frac{\pi_{\rm HMT}}{\varphi^2}.}
\]
With volumetric normalization, the relative difference of readers is
\[
 \frac{\mathcal L_{\rm ar}-\mathcal L_{\rm vol}}
 {\mathcal L_{\rm vol}}
 =
 \frac{\pi_{\rm HMT}}{\varphi^2}-1
 =
 0.1999816148643266611\ldots.
\]
The physical reading of this correspondence is undertaken later, in the
Cartan--Holst chapter, after constructing holonomy, co-tetrad, connection, torsion,
and curvature. Here the result is entirely HMT: an exact ratio between
two readers of the same sheet memory.

They should be distinguished from the enriched extractor subsequently used in
Dimensional Mechanics. There a complete history preserves the action count
\(n_A\), the twelve oriented events \(e_0,\ldots,e_{11}\) and the torsional
cochain at step level; these yield
\[
 s_C=\sum_{i=0}^{5}(e_i+e_{i+6}),
 \qquad
 W_\pm=6n_A\pm s_C.
\]
The channels already generated by HMT then satisfy the exact identity
\[
 \exp\!\left(n_AA_{\rm rad}+\frac{s_C}{6}C^\ast_{\rm rad}\right)
 =\left(q_+^{-W_+}q_-^{-W_-}\right)^{1/12}.
\]
Therefore, the route--angular-character passage is not absent: it is realized on
the enriched dodecaphase record. What this phase calculation proves by itself
are \(\Omega_{\rm sw},S,G\); it neither replaces them nor permits reconstruction of the
complete record after it has been projected.

The uniform measure on the \(104\,976\) presentations descends to atoms
of masses \(1/729\), \(1/243\) and \(1/54\), but not to \(1/468\). In the
fivefold microfiber of \(\pi\), the two probabilities are
\[
 \tau_{468}(\mathscr R_\pi)=\frac5{468},
 \qquad
 F_\#\mu_{\rm seed}(\mathscr R_\pi)=\frac7{729}.
\]
Both the averaged lift and the phase lift realize the first: they do not
push forward the uniform measure on seeds, but assign the same total
mass to each observable class. The phase lift also preserves
microscopic stillness at neutral steps.

To relate phase elevation to the enriched histories of
depth six, a kernel \(\mathcal K_6\) on
\(\mathcal H_6^{\rm cat}\times C_9\) is required. It descends to
\(\widehat{\mathcal K}_{\rm ph}\) through
\(\rho_6\times\mathrm{id}_{C_9}\) if and only if it satisfies, at each phase, the
strong lumpability condition
\[
 \sum_{\rho_6(g)=y}\mathcal K_6((h,r),(g,r+1))
 =L_{\tau(r)}(\rho_6(h),y),
\]
independently of \(h\) within each fiber of \(\rho_6\). In that case,
composition with \(F\) projects to \(\mathcal K_{\rm ph}\). Constructing
a family \((\mathcal K_m)_{m\geq6}\) that satisfies this identity and
commutes with all truncations is the remaining step toward a
transfer of the inverse limit.

\section{Scale cocycle and cylindrical conservation}
\label{sec:cociclo-parry-escala}

Let \(r=(1,\varphi)\) be a Perron--Frobenius vector of
\(F_{\rm av}\). The Parry transition matrix is
\[
P_{ij}=\frac{(F_{\rm av})_{ij}r_j}{\varphi r_i}
=
\begin{pmatrix}
0&1\\
\varphi^{-2}&\varphi^{-1}
\end{pmatrix}.
\]
For an admissible route \(\gamma:i\to j\) of length \(n\),
\[
\mathbb P(\gamma\mid i)=\frac{r_j}{\varphi^nr_i},
\qquad
\boxed{\chi_P(\gamma)=\varphi^n\frac{r_i}{r_j}.}
\]
The telescope proves
\[
\chi_P(\gamma_2\circ\gamma_1)
=\chi_P(\gamma_2)\chi_P(\gamma_1).
\]
On a closed route, \(\chi_P=\varphi^n\); its square and its inverse
generate the combinations
\[
C_n=\frac{\chi_P^2+\chi_P^{-2}}2,\qquad
S_n=\frac{\chi_P^2-\chi_P^{-2}}2,
\]
which satisfy
\[
x_{n+1}=3x_n-x_{n-1}.
\]

Since \(F_{\rm av}\) is symmetric, let
\[
\ell=\frac{r}{1+\varphi^2},\qquad \ell^{\mathsf T}r=1.
\]
The measure of a cylinder is
\[
\mu[x_0\cdots x_n]=\ell_{x_0}r_{x_n}\varphi^{-n}.
\]
Therefore,
\[
\boxed{
V_n(x)
=\frac{\mu[x_0]}{\mu[x_0\cdots x_n]}
=\varphi^n\frac{r_{x_0}}{r_{x_n}}
=\chi_P(x_0\cdots x_n).}
\]
For a typed dimension scale \(D>0\),
\[
a_n^{(D)}(x)=V_n(x)^{1/D}
\]
satisfies the cylindrical conservation law
\[
\boxed{
\bigl(a_n^{(D)}(x)\bigr)^D
\mu[x_0\cdots x_n]=\mu[x_0].}
\]

In spatial rank three and on closed returns, we write
\(a_n:=a_n^{(3)}\) and obtain
\[
\boxed{
\frac{a_9}{a_0}=\varphi^3,\qquad
\frac{a_{27}}{a_0}=\varphi^9,\qquad
\frac{a_{108}}{a_0}=\varphi^{36}.}
\]
If \(A_k:=a_{9k}\) for \(k\geq1\), then, for \(k\geq2\),
\[
A_{k+1}-2A_k+A_{k-1}
=A_{k-1}(\varphi^3-1)^2>0.
\]
The inequality proves convexity in return depth; a physical
acceleration further requires a duration chart.

The primary computation of
\(\operatorname{Sym}^2(F_{\rm av})\) in
\cref{sec:retorno-marcado-pi-phi2} shows that its spectrum is
\(\{\varphi^2,-1,\varphi^{-2}\}\). Here we conserve only the corollary
for the scale cocycle: on a closed route with \(V_n=\varphi^n\), the
stable mode obeys
\[
\boxed{M_n=\varphi^{-2n}=V_n^{-2}.}
\]
Identification of these three eigenspaces with the three TRIT sheets
requires an intertwiner; coincidence of cardinality does not replace it. The
corollary would be refuted if the closed route failed to satisfy
\(V_n=\varphi^n\) or if the stable mode did not have eigenvalue
\(\varphi^{-2}\).

\section{Dynamic prolongation theorem: global hypotheses and certified domain}

The following objects have been explicitly constructed and separated from the finite problem:
\[
 F,
 \qquad
 F_{\rm obs},
 \qquad
 \mathcal K_{\rm ph},
 \qquad
 \widehat{\mathcal K}_{\rm ph},
 \qquad
 K_{\rm cat},
 \qquad
 \tau_{468},
 \qquad
 F_{\rm av},
 \qquad
 \operatorname{Sym}^2(F_{\rm av}),
 \qquad
 D_\pi,
 \qquad
 \mathcal Q.
\]
For \(F_{\rm obs}\), deterministic return has been proved and an
explicit witness of failure to descend as a product has been exhibited. For \(\mathcal K_{\rm ph}\),
double stochasticity, irreducibility, period nine, uniform
nonadic return, and uniqueness of
\(\pi_{\rm ph}\) have been proved. For \(\widehat{\mathcal K}_{\rm ph}\), strong
lumpability, microscopic conservation of rest, and the compensated
stationary measure have been proved. \(K_{\rm cat}\) is no longer an average without
provenance: it is the shadow of one step after the uniform phase is averaged.
The TPK mode quotient induces \(F_{\rm av}\); its quadratic lift
separates the modes \(\varphi^2,-1,\varphi^{-2}\). The operator \(D_\pi\)
marks the boundary return exactly once with the generated closure, and the
interval reader \(\mathcal Q\) preserves the nesting of the cylinders of
\(\pi\) and \(\varphi\). These four objects construct
\(\pi_{\rm HMT}/\varphi_{\rm HMT}^2\) without conflating temporal frequency,
the uniform measure, and the Parry measure.

The force of this induction is relative and concrete. The TPK word forces the
weights \(1/3\). The APP interpretation of all compatible continuations,
together with uniform counting, forces \(E_+\) and \(E_\times\). The deterministic
trajectory \(F_{\rm obs}\) does not force those operators and, indeed, does not
produce them. This is not a defect of the result, but the mathematical difference
between evolution along a path and transfer over a space of paths.

\ifdefined\HMTInternalTraceability
\begin{criterion}[Enriched lift of reversion]
\label{pf84:SYN31-RHO-LIFT}
A subdomain \(D_\rho\subset X_{\rm TPK}^{\rm enr}\) and a map are sought
\[
 \Theta_\rho:D_\rho\longrightarrow X_{\rm TPK}^{\rm enr}
\]
that lifts the integer reversal \(27\mapsto72\) of
\cref{int:pf84:SYN30-RHO-DESC,int:pf84:SYN30-CROSS-WITNESS}. It must commute with the
residual, integer, and genealogical readers, quantify the quotient and the
carry, and preserve orientation, route, return, and, when the transport is
invertible, monodromy, region, and survival.
Only if \(D_\rho\) is total and invariant under \(\Theta_\rho\) may
\(\Theta_\rho^2=I\) be required and the map be called an involution.

Neither this construction nor a proof of involutivity yet exists. The
program is falsified if the map does not exist, if one of the reading squares
fails, if memory is lost without being quantified, or if
\(\Theta_\rho^2\ne I\) on a domain where involution is claimed.
\end{criterion}

\begin{criterion}[Enriched lift of the square]
\label{pf84:SYN31-SQ-LIFT}
A subdomain \(D_s\subset X_{\rm TPK}^{\rm enr}\) and a map are sought
\[
 \Theta_s:D_s\longrightarrow X_{\rm TPK}^{\rm enr}
\]
that lifts \(s(n)=n^2\), in particular \(27\mapsto729\), and commutes with the
residual, integer, and genealogical readers of
\cref{int:pf84:SYN30-SQ-DESC,int:pf84:SYN30-CROSS-WITNESS}. It must quantify
quotient and carry and preserve orientation, route, return, and, when
transport is invertible, monodromy, region, and survival. This map is not
involutive, and
\(\Theta_s^2=I\) is not imposed on it.

The proof remains absent. The nonexistence of \(\Theta_s\), the failure of
a reading square, or the uncontrolled loss of any of the
enriched coordinates would refute the program. The ablation
\(81\mapsto72\) upon changing \(q^\times(9,9)\) is an independent control:
it is a branch of neither \(\Theta_\rho\) nor \(\Theta_s\), and it proves neither
of the two maps.
\end{criterion}
\fi

The finite conclusion is complete in its domain: the uniform return, the
stationary measure, the lumpability at depth six, and the separation
between evolution of a route and transfer of the route space remain
proved by the constructed objects. In particular,
\(\widehat{\mathcal K}_{\rm ph}\) is not identified with the chronology of
\(F_{\rm obs}\), as \cref{prop:retorno-fobs-identidad} shows.
